\section{Modelos ocultos de Markov}
\label{sec:modelos_ocultos_markov}

\subsection{Modelado del tiempo}
\label{subsec:modelado_tiempo}

Asumamos que vemos el mundo de forma discreta: tomamos ``fotos'' de la
realidad cada cierto tiempo $\Tm$, conocido como \emph{tiempo de muestreo}, y
observamos lo que hay en esas fotos. Algunas cosas se pueden observar,
mientras que otras permanecen ocultas.

\begin{definicion}[Estado y evidencia]
  \label{def:estado_evidencia}
  Sea $t_k=k\,\Tm$ con $k=1,2,\dots$ Denotamos por $\X{k}$ el
  \emph{estado} del sistema en el tiempo $t_k$, el cual es un conjunto de
  variables \textbf{no observables}, y por $\Ev{k}$ la \emph{evidencia} en
  $t_k$, la cual es un conjunto de variables que \textbf{s\'i} podemos
  observar.
\end{definicion}

La notaci\'on de secuencias es la de \eqref{eq:notacion_secuencia}:
$\Xr{m}{n}=\X{m},\dots,\X{n}$ con $m\le n$. El \'indice $0$ se reserva para
el estado inicial, anterior a cualquier observaci\'on.

\subsection{Razonamiento probabil\'istico en el tiempo}
\label{subsec:razonamiento_temporal}

Cuando incorporamos el tiempo a nuestros modelos, el objeto de inter\'es es
la distribuci\'on de probabilidad de $\X{k}$ dados los estados anteriores
$\Xr{0}{k-1}$, es decir
\begin{equation}
  \label{eq:modelo_transicion_general}
  \Prob(\X{k}\given \Xr{0}{k-1}),
\end{equation}
llamada \emph{modelo de transici\'on}.

Escribir $\Xr{0}{k-1}$ dice que el estado actual depende de \emph{todos} los
anteriores: el n\'umero de padres del nodo $\X{k}$ crece sin l\'imite con
$k$, y con \'el el tama\~no de la tabla de probabilidad condicional. Un
supuesto que evita esto es pedir que $\X{k}$ dependa s\'olo de un n\'umero
finito de estados anteriores.

\begin{supuesto}[Markov de primer orden]
  \label{sup:markov}
  El estado actual depende \'unicamente del estado inmediatamente anterior:
  \begin{equation}
    \label{eq:supuesto_markov}
    \Prob(\X{k}\given \Xr{0}{k-1}) = \Prob(\X{k}\given \X{k-1}).
  \end{equation}
\end{supuesto}

\begin{supuesto}[Proceso estacionario]
  \label{sup:estacionario}
  La \emph{forma} de la distribuci\'on de probabilidad no cambia con $t_k$:
  la tabla $\Prob(\X{k}\given \X{k-1})$ es la misma para todo $k$.
\end{supuesto}

\begin{observacion}
  Los dos supuestos son independientes entre s\'i. El de Markov limita
  \emph{de cu\'antas} variables depende el estado; el de estacionariedad dice
  que los n\'umeros de esa dependencia no se reindexan con el tiempo. Juntos
  hacen que el modelo completo se especifique con una sola tabla finita,
  sin importar cu\'antos pasos de tiempo se observen.
\end{observacion}

La evidencia $\Ev{k}$ se relaciona con el estado mediante el
\emph{modelo del sensor}, para el cual se hace un supuesto an\'alogo.

\begin{supuesto}[Markov del sensor]
  \label{sup:markov_sensor}
  La evidencia en $t_k$ depende s\'olo del estado en $t_k$:
  \begin{equation}
    \label{eq:supuesto_sensor}
    \Prob(\Ev{k}\given \Xr{0}{k},\Evr{1}{k-1}) = \Prob(\Ev{k}\given \X{k}).
  \end{equation}
\end{supuesto}

A $\Prob(\Ev{k}\given \X{k})$ se le llama modelo del sensor (o modelo de
observaci\'on). N\'otese que el \'indice de la evidencia empieza en $1$: se
asume que \textbf{no hay evidencia para el tiempo $t_0$}.

\subsection{Distribuci\'on conjunta}
\label{subsec:distribucion_conjunta}

La distribuci\'on de probabilidad del estado inicial se denota
$\Prob(\X{0})$. Con ella y con los tres supuestos anteriores, la
distribuci\'on conjunta de los estados $\Xr{0}{k}$ y las evidencias
$\Evr{1}{k}$ queda completamente determinada:
\begin{equation}
  \label{eq:conjunta}
  \Prob(\Xr{0}{k},\Evr{1}{k})
  = \Prob(\X{0})\prod_{i=1}^{k}
    \Prob(\X{i}\given \X{i-1})\,\Prob(\Ev{i}\given \X{i}).
\end{equation}

\begin{observacion}
  El modelo entero se especifica entonces con tres piezas: la distribuci\'on
  inicial $\Prob(\X{0})$, el modelo de transici\'on
  $\Prob(\X{i}\given \X{i-1})$ y el modelo del sensor
  $\Prob(\Ev{i}\given \X{i})$. Los dos productos dentro de la sumatoria de
  \eqref{eq:conjunta} son justamente los que se leen de la red bayesiana
  din\'amica de la Figura~\ref{fig:red_sombrilla}.
\end{observacion}

\begin{figure*}[t]
  \centering
  \begin{tikzpicture}[
      >={Stealth[round]},
      nodo/.style={ellipse, draw=black!70, fill=blue!8, minimum width=22mm,
                   minimum height=7mm, inner sep=2pt, font=\small},
      tabla/.style={font=\scriptsize, inner sep=2pt},
      every path/.style={->, thick, draw=black!70}
    ]
    % --- estados ocultos ---
    \node[nodo] (r0) {$\mathit{Lluvia}_{k-1}$};
    \node[nodo, right=26mm of r0] (r1) {$\mathit{Lluvia}_{k}$};
    \node[nodo, right=26mm of r1] (r2) {$\mathit{Lluvia}_{k+1}$};

    % --- evidencia ---
    \node[nodo, below=14mm of r0] (u0) {$\mathit{Sombrilla}_{k-1}$};
    \node[nodo, below=14mm of r1] (u1) {$\mathit{Sombrilla}_{k}$};
    \node[nodo, below=14mm of r2] (u2) {$\mathit{Sombrilla}_{k+1}$};

    % --- transiciones ---
    \draw ($(r0)+(-16mm,0)$) -- (r0);
    \draw (r0) -- (r1);
    \draw (r1) -- (r2);
    \draw (r2) -- ($(r2)+(16mm,0)$);

    % --- sensores ---
    \draw (r0) -- (u0);
    \draw (r1) -- (u1);
    \draw (r2) -- (u2);

    % --- tablas de probabilidad condicional ---
    \node[tabla, above=5mm of r1, draw=black!40] (tt) {%
      \begin{tabular}{c|c}
        $\mathit{Ll}_{k-1}$ & $\prob(\mathit{Ll}_{k}\given \mathit{Ll}_{k-1})$\\\hline
        $\Verd$ & $0.7$\\
        $\Fal$  & $0.3$
      \end{tabular}};

    \node[tabla, right=6mm of u1, draw=black!40] (ts) {%
      \begin{tabular}{c|c}
        $\mathit{Ll}_{k}$ & $\prob(\mathit{So}_{k}\given \mathit{Ll}_{k})$\\\hline
        $\Verd$ & $0.9$\\
        $\Fal$  & $0.2$
      \end{tabular}};
  \end{tikzpicture}
  \caption{Estructura de la red bayesiana y de las distribuciones de
    probabilidad condicional del ejemplo de la sombrilla. El modelo de
    transici\'on es $\Prob(\mathit{Lluvia}_k\given \mathit{Lluvia}_{k-1})$ y
    el modelo del sensor es
    $\Prob(\mathit{Sombrilla}_k\given \mathit{Lluvia}_k)$. Adaptada
    de \cite{russell}.}
  \label{fig:red_sombrilla}
\end{figure*}

\subsection{Tipos de inferencia}
\label{subsec:tipos_inferencia}

Hay cuatro tipos de inferencias b\'asicas sobre estos modelos:

\begin{itemize}
  \item \textbf{Filtrado} (estimaci\'on de estado): obtener la
    distribuci\'on de probabilidad del estado \emph{actual} dada toda la
    evidencia hasta ahora,
    \begin{equation}
      \label{eq:filtrado}
      \Prob(\X{k}\given \Evr{1}{k}=\evr{1}{k}).
    \end{equation}

  \item \textbf{Predicci\'on}: calcular la distribuci\'on de probabilidad de
    un estado \emph{futuro} dada la evidencia disponible,
    \begin{equation}
      \label{eq:prediccion}
      \Prob(\X{k+j}\given \Evr{1}{k}=\evr{1}{k}),\qquad j>0.
    \end{equation}

  \item \textbf{Suavizado}: calcular la distribuci\'on de probabilidad de un
    estado \emph{pasado} $\X{j}$,
    \begin{equation}
      \label{eq:suavizado}
      \Prob(\X{j}\given \Evr{1}{k}=\evr{1}{k}),\qquad 0\le j\le k.
    \end{equation}

  \item \textbf{Explicaci\'on m\'as probable}: obtener la secuencia de
    estados m\'as probable que haya generado las observaciones,
    \begin{equation}
      \label{eq:explicacion}
      \argmax_{\Xr{1}{k}} \prob(\Xr{1}{k}\given \Evr{1}{k}=\evr{1}{k}).
    \end{equation}
\end{itemize}

\begin{observacion}
  Filtrado, predicci\'on y suavizado se distinguen s\'olo por la posici\'on
  del estado consultado respecto al final de la evidencia: igual, posterior o
  anterior. La explicaci\'on m\'as probable es de naturaleza distinta: no
  devuelve una distribuci\'on sino una secuencia completa, y en general
  \emph{no} coincide con encadenar los estados individualmente m\'as
  probables obtenidos por suavizado.
\end{observacion}

\subsection{Ejemplo: c\'alculo directo}
\label{subsec:ejemplo_calculo}

\begin{ejemplo}[Filtrado por enumeraci\'on en el modelo de la sombrilla]
  \label{ej:sombrilla_filtrado}
  Digamos que queremos calcular
  $\Prob(R_2\given M_{1:2}=[\Verd,\Verd])$, donde $R_k$ abrevia
  $\mathit{Lluvia}_k$ y $M_k$ abrevia $\mathit{Sombrilla}_k$, usando
  directamente la conjunta \eqref{eq:conjunta}.

  Marginalizando sobre las variables ocultas $R_0$ y $R_1$ y fijando la
  evidencia observada:
  \begin{align}
    \label{eq:sombrilla_enumeracion}
    \Prob(R_2\given M_{1:2}&=[\Verd,\Verd]) = \nonumber\\
    \al \sum_{R_1}\sum_{R_0}
      &\prob(M_2=\Verd\given R_2)\,\prob(R_2\given R_1)\nonumber\\
      &\cdot\, \prob(M_1=\Verd\given R_1)\,\prob(R_1\given R_0)\,\prob(R_0),
  \end{align}
  donde $\al$ es una constante de normalizaci\'on, elegida de modo que las
  entradas de la distribuci\'on resultante sumen $1$.
\end{ejemplo}

\begin{observacion}
  Este c\'alculo por enumeraci\'on es correcto pero no escala: el n\'umero de
  t\'erminos de la suma crece exponencialmente con $k$, pues se marginaliza
  sobre $R_0,\dots,R_{k-1}$. El algoritmo de filtrado recursivo
  (\emph{forward}) reutiliza el resultado del paso anterior y hace el trabajo
  constante por paso de tiempo.
\end{observacion}
